% =====================================================================
% Introduction
% =====================================================================
\section{Introduction}
\label{sec:introduction}

\IEEEPARstart{M}{odern} robotic systems---from industrial
manipulators on factory floors to autonomous mobile robots in
warehouses---share a fundamental computational challenge: given a
sequence of desired positions (waypoints) produced by a higher-level
path planner, how should the robot interpolate between those
positions to produce motion that is physically feasible, smooth, and
efficient? This question sits at the intersection of numerical
analysis and robotics engineering, and the choice of interpolation
method has direct consequences for mechanical wear, energy
consumption, payload stability, and whether the planned trajectory
can be executed in real time~\cite{biagiotti2008}.

The simplest answer---connecting consecutive waypoints with
straight-line segments---creates abrupt velocity changes at every
waypoint that would require infinite acceleration to follow exactly,
making the trajectory physically infeasible. More sophisticated
interpolation methods from Chapters~18, 24, and 28 of Chapra and
Canale~\cite{chapra2021} address this problem by producing curves
with guaranteed continuity in velocity, acceleration, or even jerk
(the time derivative of acceleration). Higher-order continuity
generally reduces the mechanical disturbance that the robot
experiences during motion, but every additional constraint must be
paid for somewhere, and the practical trade-offs are not always
obvious from the asymptotic analysis alone.

\subsection{Problem statement}
\label{sec:problem-statement}

A robot is given a sequence of $n+1$ waypoints
$\{(t_i, \vec{p}_i)\}_{i=0}^{n}$ where $\vec{p}_i \in \mathbb{R}^d$ is
the desired configuration at time $t_i$. We require an interpolating
function $\vec{p}(t) : [t_0, t_n] \to \mathbb{R}^d$ such that
$\vec{p}(t_i) = \vec{p}_i$ (or, in the case of approximation methods,
$\vec{p}(t_i) \approx \vec{p}_i$). The choice of $\vec{p}(\cdot)$
should be evaluated against multiple criteria: positional accuracy,
velocity continuity, acceleration smoothness, jerk minimisation,
computational cost, and total path length.

The central question this paper addresses is:

\begin{quote}
\itshape Among the numerical interpolation methods presented in
Chapters~18 and~24 of Chapra and Canale, which produces the most
suitable trajectories for robotic applications, and under what
conditions does each method excel or fail?
\end{quote}

\subsection{Contributions}
\label{sec:contributions}

This paper makes the following contributions:

\begin{enumerate}[leftmargin=*]
\item We provide textbook-formulation Python implementations of four
interpolation methods---linear (Eq.~18.2 of \cite{chapra2021}),
natural and clamped cubic splines (Eq.~18.36--18.37), $C^4$ quintic
Hermite splines, and cubic B-spline approximation via the Cox--de
Boor recursion---validated against reference implementations and
analytic test functions.
\item We construct two robotics datasets (a 12-waypoint mobile-robot
aisle and an 8-waypoint 2-DOF arm pick-and-place) and seven dataset
variants that sweep over waypoint density, geometric character, and
time spacing.
\item We evaluate each method under each variant against six
quantitative metrics (peak speed, peak acceleration, RMS jerk, total
path length, maximum waypoint error, computation time) and a seventh
that quantifies $C^1$ violations directly.
\item We validate the analytic derivative formulas against
finite-difference numerical derivatives drawn from Chapter~28 of
\cite{chapra2021}, confirming agreement to better than $10^{-3}$ in
the relevant units.
\item We identify a non-obvious result: the quintic spline's
apparently inferior smoothness under natural boundary conditions is
an artefact of the boundary condition itself, not the method;
cubic-derived clamped boundary conditions recover and exceed the
cubic spline's performance.
\end{enumerate}

\subsection{Paper organisation}

The remainder of the paper is organised as follows.
Section~\ref{sec:background} reviews the relevant numerical methods
and the robotics-specific motivation for studying them.
Section~\ref{sec:methodology} describes the implementation, the
datasets, and the evaluation metrics. Section~\ref{sec:results}
presents quantitative results from the baseline scenarios, the full
variant sweep, the derivative validation, and the boundary-condition
sensitivity study. Section~\ref{sec:discussion} interprets the
findings in the context of practical trajectory generation, and
Section~\ref{sec:conclusion} concludes with limitations and future
work.
